\documentclass[12pt]{article}

\usepackage[margin=1in]{geometry}
\usepackage{setspace}
\usepackage{hyperref}
\usepackage{csquotes}

\setstretch{1.15}

\title{\vspace{-1.2cm}
EQUITAS: Enforceable Equity for AI Healthcare Diagnostics\\
\large A Justice-Based Framework Grounded in CLD Adolescent Diagnostic Harm, Community Co-Design, Binding Governance, and Civic Literacy
}
\author{Piter Garcia\\IDAI700 (UOR)}
\date{December 2025}

\begin{document}
\maketitle

% ----------------------------
\section{Introduction (\texorpdfstring{$\sim$}{\textasciitilde}400 words)}

Diagnostic systems were not built to see some lives clearly---and the cost of that invisibility is not abstract. [file:4]
As a Latino (Dominican) immigrant and trauma survivor who later lived through chronic illness, depression, and years of missed ADHD recognition, repeated clinical encounters taught a single lesson: when institutions cannot read a person well, they often recode distress as defiance, laziness, or ``psychosomatic noise.'' [file:4]
In my case, delayed and wrong diagnosis compounded harm: symptoms were treated without attending to cultural context, linguistic friction, or trauma-informed explanations, and the result was not merely clinical error but a predictable outcome of systems optimized for a presumed ``normal'' patient. [file:4]

This paper argues that AI-driven healthcare diagnostics risk automating that same institutional misrecognition at scale unless equity is made enforceable rather than aspirational. [file:4]
Within IDAI700, the ethical stakes are best framed as a convergence of \textit{epistemic injustice} (whose knowledge counts in defining symptoms and normality), \textit{procedural justice} (who has decision authority over design and deployment), and \textit{distributive justice} (who bears harms and who receives benefits). [file:6]
The central claim is that trustworthy diagnostic AI for culturally and linguistically diverse (CLD) adolescents cannot be achieved through principles alone because principles can remain meaningless, isolated, and toothless in practice. [file:7]

To answer this, I propose \textbf{EQUITAS}, a novel framework that operationalizes enforceable equity through: (1) community co-design with decision authority, (2) binding governance with auditability and veto gates, and (3) civic literacy via Digital Resistance Mapping so communities can interpret, contest, and reshape diagnostic technologies. [file:4]
EQUITAS deliberately integrates trust-building co-design work (Nadarzynski et al.), governance infrastructure (Lekadir et al.'s FUTURE-AI), measurable trust indicators (Sagona et al.), empirical bias evidence (Marko et al.), ethics-to-engineering translation (Chen et al.), and IDAI700 skepticism about ``ethics theater'' (Munn) into a single actionable model. [file:4][file:7]
The thesis is not that AI should be rejected, but that diagnostic AI must be treated as a justice-critical institution: if it cannot demonstrate equity for CLD adolescents with intersecting identities (ADHD, chronic illness, depression, childhood trauma), it is not safe to deploy. [file:4]

% ----------------------------
\section{Background: Synthesis of Six Sources (\texorpdfstring{$\sim$}{\textasciitilde}800 words)}

\subsection*{Nadarzynski et al. (2024): Co-Design as the Origin of Trust}
Nadarzynski et al. provide the participation spine: co-design is not a late-stage ``acceptance'' exercise but the origin point of legitimacy for AI health services. [file:4]
Their lifecycle approach embeds marginalized community members across stages---from problem framing and prototyping to piloting, deployment, and retirement---so that local realities shape both system goals and harm expectations. [file:4]
For CLD adolescents, the ethical relevance is immediate: if families and youth cannot shape what counts as a symptom, a stress response, or a trauma adaptation, then the diagnostic model inherits majority-culture assumptions as ground truth. [file:4]

\subsection*{Lekadir et al. (2025): FUTURE-AI as Governance Backbone}
Lekadir et al.'s FUTURE-AI translates ``trustworthy AI'' into governance infrastructure with concrete, auditable practices. [file:4]
The framework specifies pillars and operational tools (e.g., documentation artifacts, provenance tracking, subgroup reporting, post-deployment monitoring, and retirement criteria) that move ethics from slogans to checkable obligations. [file:4]
However, as my course reflections emphasize, governance can still preserve institutional autonomy if affected communities have no binding authority to block or withdraw harmful systems. [file:4]

\subsection*{Sagona et al. (2025): Trust as an Outcome You Can Measure}
Sagona et al. treat trust not as a vibe or marketing goal but as something that can be operationalized through indicators of trustworthiness. [file:4]
This matters ethically because institutions often demand trust from marginalized communities while refusing to earn it through demonstrated performance, transparency, and recourse. [file:4]
For CLD adolescents navigating stigma, language barriers, and clinical dismissal, ``trust'' must be tethered to measurable practices: clarity about what the system does, evidence of subgroup performance, and proof of meaningful contestability. [file:4]

\subsection*{Marko et al. (2025): The Empirical Backbone}
Marko et al. provide empirical grounding: healthcare AI systems frequently underperform for marginalized groups because inclusivity is absent from data, teams, and evaluation. [file:4]
In my paper draft materials, Marko et al. are described as synthesizing a large body of studies documenting bias and disparities, reinforcing that diagnostic inequity is not anecdotal but patterned and reproducible. [file:4]
This is ethically decisive: when harms are documented, continued deployment becomes a governance choice rather than an uncertainty problem. [file:4]

\subsection*{Chen et al. (2023): Turning Ethics into Engineering}
Chen et al. focus on translating ethical intent into computational practice via curriculum and engineering workflows. [file:4]
This directly addresses a core failure mode highlighted in IDAI700: ethical language often fails at the ``handoff'' point where developers need implementable requirements, tests, and artifacts. [file:7]
Yet EQUITAS extends Chen's insight by insisting that translation is not only technical (principle-to-feature) but political (community-to-authority): whose ethics are being implemented, and who can stop deployment when those ethics are violated. [file:4]

\subsection*{Munn (IDAI700 Unit 4) + Digital Resistance Mapping: From Principles to Power}
Munn's critique---that prevailing AI ethics is often meaningless, isolated, and toothless---clarifies why many well-intentioned frameworks fail to change outcomes. [file:7]
In response, my Digital Resistance Mapping approach adds a pedagogical layer: CLD communities need civic literacy tools that help them map where harm occurs (data collection, labeling, clinical interpretation, feedback loops, procurement, enforcement) and where intervention is possible. [file:7]
This is not ``education'' as a substitute for governance; it is education as a precondition for legitimate consent and community power, aligning with course foundations that treat transparency, accountability, and trust as institutionally produced rather than individually owed. [file:6][file:7]

\subsection*{Critical synthesis}
Together, these sources align on a key insight: trustworthiness must be built through process (co-design), infrastructure (governance artifacts), evidence (bias documentation), and operationalization (engineering translation), not asserted through rhetoric. [file:4][file:7]
They also reveal a tension: participation without authority risks tokenism, while governance without enforcement risks ethics-washing. [file:4][file:7]
EQUITAS is designed to resolve that tension by making equity a deployment condition with community veto, measurable trust indicators, and civic capacity to exercise those rights. [file:4]

% ----------------------------
\section{Your Argument: The EQUITAS Framework (\texorpdfstring{$\sim$}{\textasciitilde}1000 words)}

\subsection{Part A: The Problem (\texorpdfstring{$\sim$}{\textasciitilde}300 words)}
CLD adolescents experience diagnostic disparity not simply because clinicians or tools ``make mistakes,'' but because the epistemic and procedural architecture of diagnosis was historically standardized on majority-culture norms. [file:4]
My ED415 research brief is treated in my draft as evidence that CLD adolescents are underdiagnosed for ADHD and depression despite comparable or higher symptom burden, and that screening tools perform worse for ethnic-minority and non-English-speaking youth. [file:4]
When ADHD, chronic illness, depression, and childhood trauma intersect, misrecognition compounds: trauma-informed hypervigilance can be misread as oppositionality, executive dysfunction can be moralized as laziness, and chronic symptoms can be dismissed as psychosomatic rather than investigated. [file:4]

AI does not enter a neutral space; it enters a diagnostic ecosystem already shaped by power. [file:4]
If an algorithm is trained on data that encodes majority-culture presentations as ``normal'' and CLD presentations as noise or noncompliance, then even sophisticated fairness metrics can coexist with lived invisibility. [file:4]
This is the governance illusion: an institution can publish equity language, deploy dashboards, and hold ethics meetings while preserving the underlying power structure that produces inequity. [file:4]

Munn's critique sharpens why this persists: ethics frameworks without enforcement mechanisms can become toothless, enabling institutions to claim virtue while continuing business as usual. [file:7]
In healthcare, this dynamic is intensified by utilitarian pressure to deploy ``good enough'' systems quickly, often treating marginalized error rates as acceptable collateral. [file:6]
EQUITAS begins from the opposite premise: if documented diagnostic harm concentrates in CLD adolescents, then the ethical default is precaution and non-deployment until enforceable equity conditions are met. [file:4]

\subsection{Part B: The Solution (\texorpdfstring{$\sim$}{\textasciitilde}400 words)}
EQUITAS operationalizes enforceable equity through three pillars.

\subsubsection*{Pillar 1: Equity-Centered Co-Design with Decision Authority}
Following Nadarzynski et al., EQUITAS requires co-design from problem framing through retirement, but it upgrades co-design from participation to authority. [file:4]
Practically, this means CLD adolescents and families are not merely consulted; they hold seats with binding votes on what the tool is for, what counts as harm, what language is acceptable, and what tradeoffs are forbidden. [file:4]
This directly addresses epistemic injustice by treating lived experience as primary data about symptom meaning, cultural adaptation, and trauma expression rather than as ``anecdotal feedback.'' [file:4]

\subsubsection*{Pillar 2: Binding Governance + Enforcement (CAB Veto Power)}
EQUITAS adopts FUTURE-AI as a governance backbone (documentation, traceability, subgroup performance reporting, monitoring, and retirement triggers), but adds an enforcement layer that institutions cannot unilaterally waive. [file:4]
A Community Accountability Board (CAB) has veto authority at procurement and at deployment gates, including the power to pause or retire a model when subgroup harms exceed thresholds. [file:4]
This directly answers Munn's toothless critique by converting ethics from a voluntary code into a governance mechanism with real consequences. [file:7]

\subsubsection*{Pillar 3: Civic Literacy via Digital Resistance Mapping}
Trustworthy deployment requires more than transparency; it requires that affected communities can interpret and use what is disclosed. [file:7]
Digital Resistance Mapping is the civic literacy layer: a structured method for CLD youth and families to map where algorithmic decisions come from (data sources, labels, clinical interfaces), where harm accumulates, and what recourse pathways exist. [file:7]
This enables procedural justice by making participation informed and actionable rather than symbolic, and it reduces the ``transaction costs'' of contesting a diagnostic system by clarifying who is responsible for what. [file:7]

\subsubsection*{Metrics and trust indicators}
Borrowing from Sagona et al., EQUITAS treats trust as an outcome to be measured, not demanded. [file:4]
At minimum, EQUITAS requires disaggregated performance reporting for language and cultural subgroups, documentation artifacts legible to non-experts, and post-deployment monitoring for drift that can trigger automatic review. [file:4]
Crucially, metrics become enforceable because the CAB can act on them, not merely observe them. [file:4]

\subsection{Part C: Why This Approach is Ethically Sound (\texorpdfstring{$\sim$}{\textasciitilde}300 words)}
EQUITAS is grounded in standpoint epistemology as practiced in this project: my lived experience is not treated as decoration, but as a form of situated knowledge about how diagnostic systems fail CLD adolescents. [file:4]
In IDAI700 ethics foundations, rights-based and justice-based lenses support the claim that marginalized communities have a moral right to equitable diagnostic systems and to meaningful accountability when harmed. [file:6]
Care ethics also supports EQUITAS because adolescent diagnosis is relational and context-dependent; it cannot be reduced to decontextualized outputs without violating what care requires in practice. [file:6]

EQUITAS is also structurally responsive to power analysis emphasized in the course: shifting outcomes requires shifting authority, not merely adding fairness language. [file:7]
By embedding veto power and enforceable gates, EQUITAS treats equity as a safety requirement rather than a competitive feature. [file:4]
Finally, by connecting ethics to engineering practices (via governance artifacts and measurable indicators), EQUITAS answers the implementation gap: it makes equity auditable, contestable, and stoppable when violated. [file:4][file:7]

% ----------------------------
\section{Counterargument \& Response (\texorpdfstring{$\sim$}{\textasciitilde}500 words)}

\subsection*{The pragmatic/utilitarian objection}
A strong objection (and the one explicitly anticipated in IDAI700 course guidance) is utilitarian and pragmatic: EQUITAS is too slow and expensive. [file:6]
In a resource-constrained healthcare system, the argument goes, deploying a ``good enough'' diagnostic AI quickly could help the majority, even if 15--20\% of patients (often marginalized subgroups) face worse performance. [file:6]
Delaying deployment for enforceable equity could therefore increase net harm by withholding benefits (earlier detection, triage support, reduced clinician workload) from the many in order to avoid disparities for the few. [file:6]
On this view, insisting on community veto and extensive co-design is ethically indulgent: it diverts resources from care delivery and may reduce overall welfare. [file:6]

\subsection*{Refutation (\texorpdfstring{$\sim$}{\textasciitilde}350 words)}
This utilitarian calculus is ethically distorted because it systematically undervalues marginalized lives by treating predictable subgroup harm as an acceptable externality. [file:4]
The ``15--20\%'' is not random: it tends to map onto communities already exposed to diagnostic neglect, language exclusion, and racialized patterns of dismissal, meaning the proposal repeats historical injustice under a new technical banner. [file:4]
EQUITAS rejects this because the ethical problem is not merely outcome totals but the legitimacy of procedures that allocate risk without consent and without recourse. [file:4]

Second, the pragmatic objection relies on short time horizons that ignore downstream costs of inequitable deployment: misdiagnosis and missed diagnosis create chronic cascades (delayed treatment, worsened functioning, avoidable crises) and deepen mistrust that reduces future care-seeking and research participation. [file:4]
Munn's critique helps clarify why ``principles without teeth'' are not neutral compromises: they can function as ethics-washing that delays real accountability while harms persist. [file:7]
If harms are documented and predictable, precaution is not perfectionism; it is a basic safety stance consistent with responsible governance. [file:4]

Third, EQUITAS can be cost-rational \textit{because} it reduces preventable failure: co-design catches mismatch earlier, governance artifacts reduce organizational confusion about responsibility, and measurable trustworthiness indicators reduce reputational and legal risk. [file:4][file:7]
Most importantly, EQUITAS changes who gets to decide: CLD communities become agents in defining acceptable risk and acceptable tradeoffs, rather than passive recipients of institutional cost-saving decisions. [file:4]

\subsection*{Acknowledging the tension (\texorpdfstring{$\sim$}{\textasciitilde}150 words)}
Resource scarcity is real, and some healthcare settings operate under extreme staffing and funding constraints. [file:6]
EQUITAS does not deny this; it treats scarcity as a governance problem (how resources are allocated and whose harms are tolerated) rather than as an excuse to normalize predictable injustice. [file:4]
The appropriate response is institutional investment: fund language access, community partnership infrastructure, and equity audits as core clinical safety functions, not optional add-ons. [file:4]
If a system cannot afford to deploy diagnostic AI without harming CLD adolescents, that is evidence the system cannot afford \textit{that deployment} ethically. [file:4]

% ----------------------------
\section{Conclusion (\texorpdfstring{$\sim$}{\textasciitilde}300 words)}

EQUITAS is necessary, not optional, because diagnostic AI is not merely a tool; it is an institution-like decision system that shapes who is seen, who is believed, and who receives care. [file:4]
For CLD adolescents living at the intersection of ADHD, chronic illness, depression, and childhood trauma, the moral stakes are intensified: misrecognition is not a minor error but a life-shaping redistribution of opportunity, safety, and self-understanding. [file:4]
The ethical failure at issue is therefore not only biased outputs but unjust power: who defines symptoms, who sets thresholds, and who can force change when harm occurs. [file:4][file:7]

This paper has argued that existing ethics frameworks often fail because they do not bind institutions to equity outcomes or grant communities enforceable authority. [file:7]
EQUITAS answers that failure by integrating (1) co-design as the origin of trust, (2) governance as auditable infrastructure, (3) trust as measurable outcome, (4) empirical evidence of disparities, (5) ethics-to-engineering translation, and (6) resistance mapping as civic literacy that enables communities to exercise rights in practice. [file:4][file:7]
Concretely, healthcare institutions and AI vendors should operationalize EQUITAS through procurement requirements, CAB veto gates, subgroup performance thresholds, continuous monitoring with retirement triggers, and public-facing documentation that is legible to the communities most affected. [file:4]

The call to action is not simply ``build fairer models'' but to change who gets to decide what ``fair'' means and what happens when fairness fails. [file:4]
My trajectory from diagnostic harm to framework design is not presented as inspirational narrative; it is evidence that lived experience can function as a diagnostic instrument for institutional failure and a blueprint for institutional repair. [file:4]
If AI is to have a legitimate role in healthcare diagnostics, it must be governed so that CLD adolescents are no longer rendered invisible by design. [file:4]

% ----------------------------
\section*{AI Reflection (\texorpdfstring{$\sim$}{\textasciitilde}300 words; excluded from paper word count)}

AI was used in this project primarily as a drafting and structure assistant: to transform an outline into section-level prose, to tighten the counterargument into its strongest utilitarian form, and to improve cohesion across the EQUITAS pillars (co-design, governance/enforcement, and civic literacy). [file:7]
It was also used to test whether my claims were operational rather than aspirational---for example, by forcing explicit articulation of what a CAB veto gate would mean procedurally and what kinds of measurable indicators could plausibly constitute ``trust'' as an outcome. [file:4]
In that sense, using AI mirrored the paper's theme: principles are easy to state, but operationalization requires translation into concrete requirements. [file:7]

The main limitation encountered is that AI can make justice language sound fluent while quietly smoothing over power conflicts. [file:7]
That risk is structurally similar to what Munn critiques as ethics-washing: persuasive abstraction that can appear complete while remaining unenforceable. [file:7]
To counteract that tendency, I treated my lived experience and course materials as constraints: every time the prose drifted toward generic claims, I forced it back into enforceable mechanisms (authority, veto, thresholds, monitoring, and recourse). [file:4][file:7]

This experience also changed my perspective on ``trustworthy AI'' by making the interpretability problem personal: even when AI-generated writing is coherent, it can obscure provenance and decision pathways unless deliberately documented. [file:7]
That maps directly onto why EQUITAS insists on traceability and civic literacy: communities cannot consent to or contest systems they cannot meaningfully read. [file:7]
Overall, using AI reinforced the course lesson that ethical outcomes are not produced by good intentions or smooth language; they are produced by governance, accountability, and power-sharing that survives institutional incentives. [file:7]

% ----------------------------
\section*{References (APA-style; fill with your portfolio's exact bibliographic data)}
\begin{thebibliography}{9}

\bibitem{nadarzynski2024}
Nadarzynski, T., et al. (2024).
(Study on marginalized community co-design in AI health services; full details in IDAI700 portfolio). [file:4]

\bibitem{lekadir2025}
Lekadir, K., et al. (2025).
(FUTURE-AI consensus on trustworthy AI principles and governance for medical imaging; full details in IDAI700 portfolio). [file:4]

\bibitem{sagona2025}
Sagona, A., et al. (2025).
(Quantitative trustworthiness/trust metrics for AI; full details in IDAI700 portfolio). [file:4]

\bibitem{marko2025}
Marko, N., et al. (2025).
(Systematic evidence on bias and racial disparities in clinical AI; full details in IDAI700 portfolio). [file:4]

\bibitem{chen2023}
Chen, X., et al. (2023).
(Translating ethics into computational/engineering practice through curriculum; full details in IDAI700 portfolio). [file:4]

\bibitem{munn_unit4}
Munn, L. (n.d.).
\textit{The Uselessness of AI Ethics}. IDAI700 Unit 4 course reading (PDF in course materials). [file:7]

\bibitem{garcia_ed415}
Garcia, P. (2025).
\textit{Portfolio Entry 7: The Puzzle Plan in Practice (ED415 Research Brief)}. Unpublished course research brief. [file:4]

\end{thebibliography}

\end{document}
